% =============================================================================
% PREAMBULE DU COURS NKENTSEU -- composition XeLaTeX
% -----------------------------------------------------------------------------
% XeLaTeX et non pdfLaTeX : le document est en francais accentue et utilise les
% polices systeme (Segoe UI, Consolas), ce que seul un moteur Unicode sait
% faire sans table de substitution.
% Les couleurs sont RELEVEES sur le logo Rihen, pas approximees a l'oeil :
% sarcelle #0A555F et orange #F79A28 sont les pixels exacts du fichier.
% =============================================================================

\usepackage{fontspec}
\usepackage{polyglossia}
\setdefaultlanguage{french}

% Ce cours-ci porte des FORMULES : amsmath et amssymb sont indispensables, et
% ils n'etaient necessaires a aucun des cours precedents (pas une mathematique
% ecrite). unicode-math n'est PAS charge a dessein : il remplacerait la fonte
% mathematique par une police OpenType, et les sommes, fractions et integrales
% sont mieux servies par les fontes AMS classiques.
\usepackage{amsmath}
\usepackage{amssymb}

\usepackage{geometry}
\geometry{a4paper, top=2.6cm, bottom=2.4cm, left=2.5cm, right=2.5cm,
          headheight=15pt, footskip=1.2cm}

\usepackage{xcolor}
\definecolor{rihenTeal}{RGB}{10,85,95}
\definecolor{rihenOrange}{RGB}{247,154,40}
\definecolor{rihenTealPale}{RGB}{232,241,242}
\definecolor{rihenOrangePale}{RGB}{253,243,229}
\definecolor{codeBg}{RGB}{247,248,249}
\definecolor{codeBorder}{RGB}{219,224,227}
\definecolor{codeKey}{RGB}{10,85,95}
\definecolor{codeStr}{RGB}{163,78,20}
\definecolor{codeCom}{RGB}{110,122,128}
\definecolor{codeNum}{RGB}{140,60,120}
\definecolor{grisTexte}{RGB}{58,64,68}

% Polices Windows garanties : ne pas dependre d'un paquet a installer.
\setmainfont{Segoe UI}[
  BoldFont={Segoe UI Semibold},
  ItalicFont={Segoe UI Italic},
  BoldItalicFont={Segoe UI Bold Italic}]
\setsansfont{Segoe UI}[BoldFont={Segoe UI Semibold}]
\setmonofont{Consolas}[Scale=0.92]

\usepackage{titlesec}
\usepackage{titletoc}
\usepackage{fancyhdr}
\usepackage{graphicx}
\usepackage{tikz}
\usetikzlibrary{calc}
\usepackage{listings}
\usepackage{tcolorbox}
% `listings` est indispensable ici : c'est elle qui fournit \newtcblisting,
% et donc les blocs de code encadres et coupables en fin de page.
\tcbuselibrary{skins, breakable, listings}
% Les images sont cherchees ici, quel que soit le repertoire de compilation :
% un chemin relatif casse des qu'on change -output-directory.
\graphicspath{{../assets/}{../../assets/}{./}}
\usepackage{enumitem}
\usepackage{booktabs}
\usepackage{longtable}
\usepackage{array}
\usepackage{microtype}
\usepackage[hidelinks]{hyperref}
\hypersetup{
  colorlinks=true, linkcolor=rihenTeal, urlcolor=rihenTeal,
  citecolor=rihenTeal, pdfborder={0 0 0}}
% xurl (APRES hyperref) : autorise la coupure des chemins de fichiers a
% n'importe quel caractere. Sans lui, un chemin de depot long est un bloc
% insecable qui deborde de 5 a 7 cm dans la marge.
\usepackage{xurl}
% Filet de securite pour les paragraphes charges d'identifiants en fonte a
% chasse fixe : autorise un etirement d'urgence plutot qu'un debordement.
\emergencystretch=2.5em

% ── TITRES ───────────────────────────────────────────────────────────────────
% Le numero de chapitre est POSE A COTE du titre, en gros et en pale : il sert
% de repere quand on feuillette, sans voler la vedette au titre lui-meme.
\titleformat{\chapter}[display]
  {\normalfont\sffamily\bfseries\color{rihenTeal}}
  {\hfill\fontsize{58}{58}\selectfont\color{rihenTealPale}\thechapter}
  {-14pt}
  {\fontsize{24}{28}\selectfont\raggedright}
  [\vspace{4pt}{\color{rihenOrange}\titlerule[2.4pt]}]
\titlespacing*{\chapter}{0pt}{-24pt}{28pt}

\titleformat{\section}
  {\normalfont\sffamily\bfseries\large\color{rihenTeal}}{\thesection}{0.8em}{}
\titleformat{\subsection}
  {\normalfont\sffamily\bfseries\color{grisTexte}}{\thesubsection}{0.7em}{}
\titleformat{\subsubsection}
  {\normalfont\sffamily\itshape\color{grisTexte}}{}{0pt}{}

% ── EN-TETES ET PIEDS ────────────────────────────────────────────────────────
\pagestyle{fancy}
\fancyhf{}
\fancyhead[L]{\footnotesize\sffamily\color{rihenTeal}Nkentseu — NkCanvas \& NKGui}
\fancyhead[R]{\footnotesize\sffamily\color{grisTexte}\leftmark}
\fancyfoot[C]{\footnotesize\sffamily\color{grisTexte}\thepage}
\renewcommand{\headrulewidth}{0.4pt}
\renewcommand{\headrule}{\hbox to\headwidth{\color{codeBorder}\leaders\hrule height \headrulewidth\hfill}}
\fancypagestyle{plain}{\fancyhf{}\fancyfoot[C]{\footnotesize\sffamily\color{grisTexte}\thepage}%
  \renewcommand{\headrulewidth}{0pt}}

% ── CODE ─────────────────────────────────────────────────────────────────────
% `literate` est INDISPENSABLE : sans lui, listings sous XeLaTeX avale les
% accents des commentaires francais du depot -- et tous nos exemples en ont.
\lstdefinestyle{nkcpp}{
  language=C++,
  basicstyle=\ttfamily\footnotesize\color{grisTexte},
  keywordstyle=\color{codeKey}\bfseries,
  stringstyle=\color{codeStr},
  commentstyle=\color{codeCom}\itshape,
  numberstyle=\ttfamily\scriptsize\color{codeCom},
  numbers=left, numbersep=9pt, stepnumber=1,
  showstringspaces=false,
  breaklines=true, breakatwhitespace=true,
  breakindent=14pt,
  tabsize=4,
  columns=fullflexible, keepspaces=true,
  extendedchars=true,
  frame=none, xleftmargin=13pt, framexleftmargin=13pt,
  morekeywords={uint8,uint16,uint32,uint64,int8,int16,int32,int64,float32,
                float64,usize,bool32,NkString,NkVector,NkImage,NkRect,NkColor,
                nullptr,override,constexpr,noexcept},
  literate=
    {é}{{\'e}}1 {è}{{\`e}}1 {ê}{{\^e}}1 {ë}{{\"e}}1
    {à}{{\`a}}1 {â}{{\^a}}1 {ä}{{\"a}}1
    {î}{{\^i}}1 {ï}{{\"i}}1
    {ô}{{\^o}}1 {ö}{{\"o}}1
    {ù}{{\`u}}1 {û}{{\^u}}1 {ü}{{\"u}}1
    {ç}{{\c c}}1 {É}{{\'E}}1 {È}{{\`E}}1 {Ê}{{\^E}}1
    {À}{{\`A}}1 {Ç}{{\c C}}1
    {«}{{\guillemotleft}}1 {»}{{\guillemotright}}1
    {–}{{--}}1 {—}{{---}}1 {’}{{'}}1 {…}{{\ldots}}1
}
\lstset{style=nkcpp}

% Bloc de code encadre, avec un titre qui dit D'OU vient l'extrait. Un exemple
% sans provenance ne peut pas etre verifie par le lecteur.
\newtcblisting{nkcode}[2][]{
  listing only, listing style=nkcpp,
  breakable, enhanced,
  colback=codeBg, colframe=codeBorder, boxrule=0.6pt, arc=2pt,
  left=2pt, right=2pt, top=4pt, bottom=4pt,
  % \NoAutoSpacing : le francais insere une espace fine avant « : », ce qui
  % ecrivait « fichier.cpp :412 » dans les titres de code. Une reference de
  % ligne n'est pas de la ponctuation francaise.
  title={\sffamily\bfseries\footnotesize\color{rihenTeal}\NoAutoSpacing #2},
  coltitle=rihenTeal, colbacktitle=rihenTealPale,
  fonttitle=\footnotesize, titlerule=0pt,
  attach boxed title to top left={xshift=6pt, yshift=-2pt},
  boxed title style={colframe=codeBorder, arc=1pt, boxrule=0.4pt},
  #1}

% ── ENCADRES PEDAGOGIQUES ────────────────────────────────────────────────────
% Trois seulement, et chacun a un role NET : ce qu'il faut retenir, le piege
% qui coute une soiree, et l'exercice. Multiplier les couleurs d'encadres rend
% un cours illisible -- on cesse de savoir lequel compte.
\newtcolorbox{nkretenir}[1][Ce qu'il faut retenir]{
  breakable, enhanced, colback=rihenTealPale, colframe=rihenTeal,
  boxrule=0pt, leftrule=3pt, arc=1pt,
  title={\sffamily\bfseries\small #1}, coltitle=rihenTeal,
  colbacktitle=rihenTealPale, titlerule=0pt, fonttitle=\small}

\newtcolorbox{nkpiege}[1][Le piège]{
  breakable, enhanced, colback=rihenOrangePale, colframe=rihenOrange,
  boxrule=0pt, leftrule=3pt, arc=1pt,
  title={\sffamily\bfseries\small #1}, coltitle=rihenOrange!75!black,
  colbacktitle=rihenOrangePale, titlerule=0pt, fonttitle=\small}

\newtcolorbox{nkexo}[1][Exercice]{
  breakable, enhanced, colback=white, colframe=codeBorder,
  boxrule=0.7pt, arc=2pt,
  title={\sffamily\bfseries\small #1}, coltitle=grisTexte,
  colbacktitle=codeBg, titlerule=0pt, fonttitle=\small}

% ── TABLE DES MATIERES ───────────────────────────────────────────────────────
\titlecontents{chapter}[0pt]
  {\addvspace{9pt}\sffamily\bfseries\color{rihenTeal}}
  {\contentslabel[\thecontentslabel]{1.4em}}{}
  {\titlerule*[0.7pc]{.}\contentspage}
\titlecontents{section}[1.9em]
  {\addvspace{2pt}\sffamily\small\color{grisTexte}}
  {\contentslabel{2.1em}}{}
  {\titlerule*[0.7pc]{.}\small\contentspage}

\setlist[itemize]{leftmargin=1.3em, itemsep=2pt, topsep=4pt}
\setlist[enumerate]{leftmargin=1.5em, itemsep=2pt, topsep=4pt}
\renewcommand{\arraystretch}{1.25}

% Chemin de fichier dans le corps du texte : la provenance doit se distinguer
% du code sans crier. Compose via \nkfileurl (url/xurl) : le chemin peut se
% COUPER en fin de ligne au lieu de deborder dans la marge.
\DeclareUrlCommand{\nkfileurl}{\urlstyle{tt}}
\newcommand{\nkfile}[1]{{\small\color{rihenTeal}\nkfileurl{#1}}}

% LE CODE EN LIGNE N'EST PAS DU FRANCAIS. La typographie francaise colle une
% espace fine INSECABLE avant « ; : ! ? » -- juste pour du texte, desastreux
% dans \texttt : « input.wheel = 0.f; input.wheelH = 0.f; » devenait UN bloc
% incassable qui debordait de 3 cm dans la marge. \NoAutoSpacing dans \texttt
% rend au code sa ponctuation exacte ET ses points de coupure.
% Robuste, pas fragile : \NoAutoSpacing developpe dans un TITRE de section
% s'ecrit tel quel dans la table des matieres et y explose. La version robuste
% traverse les arguments mobiles sans se developper.
\let\nkOldTexttt\texttt
\DeclareRobustCommand{\nkttCode}[1]{{\NoAutoSpacing\nkOldTexttt{#1}}}
\renewcommand{\texttt}{\nkttCode}
